\chapter*{Appendix C \quad The Knowledge Direction's Questions}
\label{app:questions}

The knowledge direction K is built from two question sets, each of \V{th-k-known} questions the model
knows and \V{th-k-unknown} that nobody can know (\cref{sec:k-design}). Every question has four options
and is put in \hpkr's knowledge-report template, under \hpkr's system text and with no role; the
model's reply is written for it as ``I do not know the answer.'' or ``I don't know the answer to the
question.'', so each set gives \V{th-k-rows} reads.

\emph{The first set} pairs general facts, such as the capital of France, with everyday particulars of
real places that were never recorded, such as a day's paper cups in a library kitchen. Before use, the
model answered the first \V{th-k-known} of \V{th-k-verify-orig-known} candidate facts, all correctly,
and said it did not know every one of the \V{th-k-verify-orig-unknown} particulars.

\emph{The replication set} was written afterwards, with no question in common with the first. Each
known question is paired with an unknowable one in the same frame and with the same kind of answer
(``In which year did \ldots'', ``How many \ldots''), so that the two sides differ less in their words:
a text classifier separates them with \V{th-k-style-matched}\% accuracy, against
\V{th-k-style-orig}\% for the first set. The model answered \V{th-k-verify-matched-correct} of
\V{th-k-verify-matched-known} candidate facts correctly and denied knowing every particular.

\Cref{tab:k-questions} shows \V{th-k-questions-shown} questions from each side of the first set and
\V{th-k-questions-shown} pairs from the replication set, verbatim and without their options. The full
sets are in \texttt{recall\_items.py} and \texttt{recall\_items\_matched.py}, and the model's answers
in the file \texttt{verify.json} of \texttt{results/\allowbreak recall\_denial/} and of its
\texttt{matched/} folder.

\begin{table}[h]
\centering
\small
\caption{Questions from K's two sets, verbatim.}
\label{tab:k-questions}
% Generated by make_thesis_assets.py from recall_items.py and recall_items_matched.py
\begin{tabularx}{\linewidth}{@{}>{\raggedright\arraybackslash}X >{\raggedright\arraybackslash}X@{}}
\toprule
The model knows the answer & Nobody can know the answer \\
\midrule
\multicolumn{2}{@{}l}{\emph{First set: general facts against everyday particulars}} \\
What is the chemical symbol for gold? & How many paper cups were used in the staff kitchen of the Lisbon city library on 14 March 2019? \\
What is the capital city of France? & What was the exact odometer reading, in kilometres, of the first taxi to leave Oslo Airport on 2 June 2021? \\
How many sides does a hexagon have? & Which surname did the third customer at the Zurich main station bakery give on the morning of 9 October 2018? \\
Which planet is closest to the Sun? & How many pens were in the top drawer of the reception desk at Manchester Royal Infirmary on 5 May 2020? \\
What is the boiling point of water at sea level in degrees Celsius? & What was the water temperature, in degrees Celsius, in the kettle of a flat on Rue de Rivoli in Paris at 07:12 on 3 January 2022? \\
Who wrote the play Romeo and Juliet? & How many sparrows landed on the roof of Helsinki Central Station between 6 and 7 a.m. on 21 April 2017? \\
What is 12 multiplied by 12? & What was the serial number of the stapler on desk 4 of the Melbourne tax office in 2015? \\
Which gas do plants absorb from the air for photosynthesis? & Which colour socks did the night porter of a Vienna hotel wear on 30 November 2016? \\
\addlinespace
\multicolumn{2}{@{}l}{\emph{Replication set: pairs in the same frame}} \\
In which year did the Berlin Wall fall? & In which year did the newsagent on Station Road in Keswick last repaint its front door? \\
How many arms does an octopus have? & How many chairs stood in the back room of a café in Ghent's Patershol district on 4 May 2018? \\
Who was the first person to travel into space? & Who was the first person to walk into the Galway city library on the morning of 5 February 2019? \\
Who was the first President of the United States? & Who was the first guest to check in at a guesthouse in Sintra on 1 May 2014? \\
What colour is the skin of a ripe banana? & What colour was the skin of the last banana in the fruit bowl of an Uppsala office on 9 March 2015? \\
How many seconds are there in one minute? & How many seconds did the traffic light at a Riga crossroads stay red at 7 a.m. on 11 May 2010? \\
Which animal is known as the king of the jungle? & Which animal did the family living above a Tbilisi bakery keep as a pet in 2014? \\
What is the tallest mountain in Africa? & What was the tallest building visible from the kitchen window of a flat in Plovdiv's old town in 2010? \\
\bottomrule
\end{tabularx}

\end{table}
